\section{Basics}
\lstset{ % Allgemeine Einstellungen für alle Codeblöcke
    language=C,% Sprache für Syntax-Highlighting
    basicstyle=\ttfamily, % Grundlegender Schriftstil
    keywordstyle=\color{blue},     % Farbe der Schlüsselwörter
    commentstyle=\color{gray},     % Farbe der Kommentare
    stringstyle=\color{red},       % Farbe der Zeichenketten
    numbers=left,                  % Zeilennummern links anzeigen
    numberstyle=\tiny\color{gray}, % Stil der Zeilennummern
    frame=single,                  % Rahmen um den Code
    breaklines=true,               % Automatischer Zeilenumbruch
    captionpos=b, % Bildunterschrift unterhalb des Codes
}

\begin{frame}[fragile] %fragile tag benötigt, da code Listings auftreten
\frametitle{Kompilieren}

\begin{itemize}
\item Kompilieren
    \begin{itemize}
        \item Quellcode wird in Bytecode übersetzt
        \item Überprüfung, ob Regeln der Sprache eingehalten wurden
    \end{itemize}
\item Java Programm kompilieren und ausführen
\begin{itemize}
        \item java Programmname.java
    \end{itemize} 
\end{itemize}
    
\end{frame}

\subsection{Variablen}

\begin{frame}[fragile] %fragile tag benötigt, da code Listings auftreten
\frametitle{Variablendeklaration}

In Java benutzen wir häufig Variablen, um Daten zu speichern.
Um eine Variable anzulegen benötigen wir einen Datentypen und einen Namen. Das Anlegen einer Variable heißt Deklaration:
\begin{lstlisting}
    int f;
    float g;
    String msg;
\end{lstlisting}
primitive Variablen, die nur deklariert wurden, kriegen einen Standardwert zugewiesen. Bei \textit{int} ist dieser beispielsweise 0
    
\end{frame}

\begin{frame}[fragile]
\frametitle{Variableninitialisierung}

Bei der Initialisierung von Variablen weisen wir diesen nun einen Wert zu:
\begin{lstlisting}
    int f = -2;
    float g = 42.69;
    String msg = "Hello World";
\end{lstlisting}

Referenzen gelten hier auch als Werte, die wir zuweisen können:
\begin{lstlisting}
    String importantMessage = msg + "!";
\end{lstlisting}
    
\end{frame}

\subsection{Datentypen}

\begin{frame}{Primitive Datentypen}
Welche primitiven Datentypen gibt es?
\begin{itemize}
\item boolean
\item byte
\item char
\item short
\item int
\item long
\item float
\item double
\end{itemize}
\end{frame}

\begin{frame}{Nicht primitive Datentypen}
\begin{itemize}
    \item Objektreferenzierung
    \begin{itemize}
        \item speichert anstelle des Objektes nur die Referenz auf dem Stack
        \item erlaubt die Einschränkung von Zugriff auf Attribute und Variablen des Objektes
        \item ist häufig nötig, da Stack Platz sehr begrenzt
    \end{itemize}
\end{itemize}

\begin{itemize}
    \item Beispiele
    \begin{itemize}
        \item String
        \item List 
        \item Arrays
        \item eigene Klassen
        \item ...
    \end{itemize}
\end{itemize}
\end{frame}

\begin{frame}[fragile]
\frametitle{Die String Klasse}
\begin{itemize}
    \item ein String Objekt enthält eine Zeichenkette in festgelegter Reihenfolge
    \item die Klasse String verfügt über Methoden, die auf dem Objekt selber aufgerufen werden können
\end{itemize}
\begin{lstlisting}
    String msg = "dies ist ein String!";
    String msg2 = msg.toUpperCase();
\end{lstlisting}
\end{frame}

\subsection{Arrays}
\begin{frame}[fragile]
\frametitle{Arrays}
Was ist ein Array?
\begin{itemize}
    \item speichert eine genaue Anzahl an Variablen desselben Datentypes
    \item hat eine feste Größe (nicht dynamisch erweiterbar)
    \item nicht initialiserte Felder kriegen einen Standardwert oder null zugewiesen
    \item indices für schnellen Zugriff
    \item wird von Java wie ein Objekt behandelt $\rightarrow$ liegt auf dem Heap
\end{itemize}
\begin{lstlisting}
    int[] array1 = new int[5];
    String[] array2 = {"A","R","R","A","Y"};
    System.out.println(array2[0]); // gibt "A" aus
    array1[3] = 42; // array1 = [0,0,0,42,0]
\end{lstlisting}
\end{frame}

\begin{frame}[fragile]
\frametitle{Mehrdimenionale Arrays}
Arrays können auch mehrdimensional sein. Das bedeutet, dass wir Arrays von Arrays haben können. Ein zweidimensionales Array könnte beispielsweise so aussehen:
\begin{lstlisting}
    int[][] matrix = {{1, 2},{3, 4},{5, 6}};
    System.out.println(matrix[2][0]); // gibt 5 aus

\end{lstlisting}
\end{frame}
